% GENERATED FILE. Edit canonical lessons, then run npm run generate:books.
% canonical-source-hash: fnv1a64:1e13ed3ff01bbadf
% canonical-lessons: KA-C465-pattana, KA-C465-rajadhani, KA-C465-pradesa, KA-C465-sthala, KA-C465-kendra

\chapter{Town and Capital}
\label{ch:ka-town-and-capital}
\hlchaptermodality{465}

\begin{chapteropening}
\textbf{By the end of this chapter:} \emph{I can say a small town, a capital city, a region, a place and a centre.}

Complete the last lesson of chapter 465: I can say a small town, a capital city, a region, a place and a centre.
\end{chapteropening}

\section[\texorpdfstring{\kn{ಪಟ್ಟಣ} (paṭṭaṇa)}{ಪಟ್ಟಣ (paṭṭaṇa)}]{\kn{ಪಟ್ಟಣ} (paṭṭaṇa) — a small town; a town}

\label{lesson:KA-C465-pattana}



\begin{quote}
Before the new one: say the Kannada for a gratuity, then the Kannada for to order.
\end{quote}

\subsection*{You'll want to know: \kn{ಪಟ್ಟಣ}}
\textbf{\kn{ಪಟ್ಟಣ}} — \emph{paṭṭaṇa} — \textquotedblleft{}a small town; a town\textquotedblright{}.

\begin{quote}
\kn{ನನ್ನ} \kn{ಊರು} \kn{ಒಂದು} \kn{ಚಿಕ್ಕ} \kn{ಪಟ್ಟಣ}.
\end{quote}

\begin{quote}
\emph{nanna ūru ondu cikka paṭṭaṇa.}
\end{quote}

\begin{quote}
My home town is a small town.
\end{quote}

Plural: \kn{ಪಟ್ಟಣಗಳು}, \emph{paṭṭaṇagaḷu}.

Bigger than a \kn{ಹಳ್ಳಿ}, smaller than a \kn{ನಗರ}.

\begin{cousinweb}
\textbf{\kn{ಪಟ್ಟಣ}} (\emph{paṭṭaṇa}) comes from Sanskrit \emph{paṭṭana}, \textquotedblleft{}a town\textquotedblright{} (CDIAL 7705); Turner notes that the Sanskrit word itself was probably taken from a Dravidian source.
\end{cousinweb}

\begin{grammarlens}[title={What you've built}]
One more word for getting around.
\end{grammarlens}

\subsection*{Your turn}
\begin{itemize}
\raggedright
  \item \emph{Say it:} \emph{paṭṭaṇa}
  \item \emph{Say it:} \emph{paṭṭaṇa}, once more
  \item \emph{Say it:} \emph{nanna ūru ondu cikka paṭṭaṇa.} — My home town is a small town.
  \item \emph{Recall:} say the Kannada for a gratuity, then the Kannada for to order, then say \emph{paṭṭaṇa} again
\end{itemize}

\begin{tcolorbox}[breakable,colback=teal!4,colframe=teal!35!black,title={Before you move on}]
What does \kn{ಪಟ್ಟಣ} mean? (\textquotedblleft{}A small town; a town\textquotedblright{}.) Say it once more.
\end{tcolorbox}
\section[\texorpdfstring{\kn{ರಾಜಧಾನಿ} (rājadhāni)}{ರಾಜಧಾನಿ (rājadhāni)}]{\kn{ರಾಜಧಾನಿ} (rājadhāni) — a capital city}

\label{lesson:KA-C465-rajadhani}



\begin{quote}
Before the new one: say the Kannada for to order, then the Kannada for a small town.
\end{quote}

\subsection*{You'll want to know: \kn{ರಾಜಧಾನಿ}}
\textbf{\kn{ರಾಜಧಾನಿ}} — \emph{rājadhāni} — \textquotedblleft{}a capital city\textquotedblright{}.

\begin{quote}
\kn{ಪ್ರತಿ} \kn{ರಾಜ್ಯಕ್ಕೆ} \kn{ಒಂದು} \kn{ರಾಜಧಾನಿ} \kn{ಇದೆ}.
\end{quote}

\begin{quote}
\emph{prati rājyakke ondu rājadhāni ide.}
\end{quote}

\begin{quote}
Every state has a capital.
\end{quote}

\kn{ರಾಜ}, a king, plus \kn{ಧಾನಿ}, a seat: where the ruler sits.

\begin{cousinweb}
\textbf{\kn{ರಾಜಧಾನಿ}} (\emph{rājadhāni}) is borrowed from Sanskrit \emph{rājadhānī}, \textquotedblleft{}a king's seat, a capital\textquotedblright{} (Monier-Williams).
\end{cousinweb}

\begin{grammarlens}[title={What you've built}]
One more word for getting around.
\end{grammarlens}

\subsection*{Your turn}
\begin{itemize}
\raggedright
  \item \emph{Say it:} \emph{rājadhāni}
  \item \emph{Say it:} \emph{rājadhāni}, once more
  \item \emph{Say it:} \emph{prati rājyakke ondu rājadhāni ide.} — Every state has a capital.
  \item \emph{Recall:} say the Kannada for to order, then the Kannada for a small town, then say \emph{rājadhāni} again
\end{itemize}

\begin{tcolorbox}[breakable,colback=teal!4,colframe=teal!35!black,title={Before you move on}]
What does \kn{ರಾಜಧಾನಿ} mean? (\textquotedblleft{}A capital city\textquotedblright{}.) Say it once more.
\end{tcolorbox}
\section[\texorpdfstring{\kn{ಪ್ರದೇಶ} (pradēśa)}{ಪ್ರದೇಶ (pradēśa)}]{\kn{ಪ್ರದೇಶ} (pradēśa) — a region; an area}

\label{lesson:KA-C465-pradesa}



\begin{quote}
Before the new one: say the Kannada for a small town, then the Kannada for a capital city.
\end{quote}

\subsection*{You'll want to know: \kn{ಪ್ರದೇಶ}}
\textbf{\kn{ಪ್ರದೇಶ}} — \emph{pradēśa} — \textquotedblleft{}a region; an area\textquotedblright{}.

\begin{quote}
\kn{ಈ} \kn{ಪ್ರದೇಶದಲ್ಲಿ} \kn{ತುಂಬಾ} \kn{ಕಾಡು} \kn{ಇದೆ}.
\end{quote}

\begin{quote}
\emph{ī pradēśadalli tumbā kāḍu ide.}
\end{quote}

\begin{quote}
There is a lot of forest in this region.
\end{quote}

\kn{ಗ್ರಾಮೀಣ} \kn{ಪ್ರದೇಶ}, a rural area.

\begin{cousinweb}
\textbf{\kn{ಪ್ರದೇಶ}} (\emph{pradēśa}) is borrowed from Sanskrit \emph{pradēśa}, \textquotedblleft{}a region, a country\textquotedblright{} (CDIAL 8668), built on \emph{dēśa}, the source of \kn{ದೇಶ}.
\end{cousinweb}

\begin{grammarlens}[title={What you've built}]
One more word for getting around.
\end{grammarlens}

\subsection*{Your turn}
\begin{itemize}
\raggedright
  \item \emph{Say it:} \emph{pradēśa}
  \item \emph{Say it:} \emph{pradēśa}, once more
  \item \emph{Say it:} \emph{ī pradēśadalli tumbā kāḍu ide.} — There is a lot of forest in this region.
  \item \emph{Recall:} say the Kannada for a small town, then the Kannada for a capital city, then say \emph{pradēśa} again
\end{itemize}

\begin{tcolorbox}[breakable,colback=teal!4,colframe=teal!35!black,title={Before you move on}]
What does \kn{ಪ್ರದೇಶ} mean? (\textquotedblleft{}A region; an area\textquotedblright{}.) Say it once more.
\end{tcolorbox}
\section[\texorpdfstring{\kn{ಸ್ಥಳ} (sthaḷa)}{ಸ್ಥಳ (sthaḷa)}]{\kn{ಸ್ಥಳ} (sthaḷa) — a place; a site}

\label{lesson:KA-C465-sthala}



\begin{quote}
Before the new one: say the Kannada for a capital city, then the Kannada for a region.
\end{quote}

\subsection*{You'll want to know: \kn{ಸ್ಥಳ}}
\textbf{\kn{ಸ್ಥಳ}} — \emph{sthaḷa} — \textquotedblleft{}a place; a site\textquotedblright{}.

\begin{quote}
\kn{ಊಟಕ್ಕೆ} \kn{ನಾವು} \kn{ಒಂದು} \kn{ಶಾಂತ} \kn{ಸ್ಥಳ} \kn{ಹುಡುಕಿದೆವು}.
\end{quote}

\begin{quote}
\emph{ūṭakke nāvu ondu śānta sthaḷa huḍukidevu.}
\end{quote}

\begin{quote}
We looked for a quiet place to eat.
\end{quote}

\kn{ಸ್ಥಳೀಯ}, local, is built on it.

\begin{cousinweb}
\textbf{\kn{ಸ್ಥಳ}} (\emph{sthaḷa}) is borrowed from Sanskrit \emph{sthala}, \textquotedblleft{}dry land, ground, a place\textquotedblright{} (CDIAL 13744); Kannada writes its \emph{l} as \kn{ಳ}.
\end{cousinweb}

\begin{grammarlens}[title={What you've built}]
One more word for getting around.
\end{grammarlens}

\subsection*{Your turn}
\begin{itemize}
\raggedright
  \item \emph{Say it:} \emph{sthaḷa}
  \item \emph{Say it:} \emph{sthaḷa}, once more
  \item \emph{Say it:} \emph{ūṭakke nāvu ondu śānta sthaḷa huḍukidevu.} — We looked for a quiet place to eat.
  \item \emph{Recall:} say the Kannada for a capital city, then the Kannada for a region, then say \emph{sthaḷa} again
\end{itemize}

\begin{tcolorbox}[breakable,colback=teal!4,colframe=teal!35!black,title={Before you move on}]
What does \kn{ಸ್ಥಳ} mean? (\textquotedblleft{}A place; a site\textquotedblright{}.) Say it once more.
\end{tcolorbox}
\section[\texorpdfstring{\kn{ಕೇಂದ್ರ} (kēndra)}{ಕೇಂದ್ರ (kēndra)}]{\kn{ಕೇಂದ್ರ} (kēndra) — a centre; a hub}

\label{lesson:KA-C465-kendra}



\begin{quote}
Before the new one: say the Kannada for a region, then the Kannada for a place.
\end{quote}

\subsection*{You'll want to know: \kn{ಕೇಂದ್ರ}}
\textbf{\kn{ಕೇಂದ್ರ}} — \emph{kēndra} — \textquotedblleft{}a centre; a hub\textquotedblright{}.

\begin{quote}
\kn{ಬಸ್} \kn{ನಿಲ್ದಾಣ} \kn{ನಗರದ} \kn{ಕೇಂದ್ರದಲ್ಲಿ} \kn{ಇದೆ}.
\end{quote}

\begin{quote}
\emph{bas nildāṇa nagarada kēndradalli ide.}
\end{quote}

\begin{quote}
The bus station is in the centre of the city.
\end{quote}

\kn{ಮಾಹಿತಿ} \kn{ಕೇಂದ್ರ}, an information centre.

\begin{cousinweb}
\textbf{\kn{ಕೇಂದ್ರ}} (\emph{kēndra}) is borrowed from Sanskrit \emph{kendra}, \textquotedblleft{}the centre of a circle\textquotedblright{}, which Sanskrit itself took from Greek \emph{kentron} (Monier-Williams).
\end{cousinweb}

\begin{grammarlens}[title={What you've built}]
One more word for getting around.
\end{grammarlens}

\subsection*{Your turn}
\begin{itemize}
\raggedright
  \item \emph{Say it:} \emph{kēndra}
  \item \emph{Say it:} \emph{kēndra}, once more
  \item \emph{Say it:} \emph{bas nildāṇa nagarada kēndradalli ide.} — The bus station is in the centre of the city.
  \item \emph{Recall:} say the Kannada for a region, then the Kannada for a place, then say \emph{kēndra} again
\end{itemize}

\begin{tcolorbox}[breakable,colback=teal!4,colframe=teal!35!black,title={Before you move on}]
What does \kn{ಕೇಂದ್ರ} mean? (\textquotedblleft{}A centre; a hub\textquotedblright{}.) Say it once more.
\end{tcolorbox}
